\chapter*{Conferencia V}
\addcontentsline{toc}{chapter}{LECTURA I - El timbre y el carácter del violín}
SEÑORES: El efecto de la edad en el timbre del violín es un tema repleto de interés. Entre todos los instrumentos musicales, en lo que respecta a la mejora del timbre con la edad y el uso, el violín destaca por encima de todos.
Se mantiene aislado. Curiosamente, la mejora del sonido en el violín va acompañada de la pérdida de valor de la madera al ser utilizada para otros fines. Debido a esta constante disminución del valor de la madera, la mejora ilimitada del sonido del violín se vuelve imposible. El hecho de la pérdida de valor de la madera confiere un lado patético a la historia del violín. Los mejores violines han sido los primeros en sucumbir ante esos enemigos insaciables: el calor y la humedad. La inexorable ley de la desintegración no ha podido detener a los mejores Maggini, ni a los mejores Amati, ni a los mejores Guarneri, ni a los mejores Stradivari. La mejor caja de resonancia, productora de un sonido excepcional, cede pronto ante los ataques de sus enemigos. Pronto, los mejores violines antiguos solo serán recordados.
Las incontables horas de placer que esos violines desgastados han brindado a la humanidad son algo indescriptible con cifras. Menos mal que tal placer no puede condensarse en una hora. Tal dulzura comprimida podría despoblar el mundo. Es una peculiaridad de la humanidad amar más aquello que proporciona mayor placer. De entre todos los objetos inanimados que contribuyen al placer humano, el violín encabeza la lista sin rival. Desde el interior de la cabaña en el rancho del Oeste hasta bajo las cúpulas doradas del Zar, el violín irradia luz.
En lo más profundo del afecto humano, el violín encuentra un amplio espacio. ¿Debe este objeto invaluable seguir siendo un sacrificio a la ley de la desintegración? ¿No debería la humanidad esforzarse por encontrar una defensa segura para el violín contra esta ley?
A la luz de mi experiencia en la búsqueda de medios de tal defensa, digo: "Es una tontería sentarse
"Se agachan y lloran sin hacer nada para salvar el violín." A la luz de mi experiencia, le digo a todo hombre que, sin experiencia, condena todos los esfuerzos por preservar las superficies internas del violín de los inevitables "estragos de la desintegración": "Están cometiendo un acto de violencia criminal contra la música". A aquellos que han hecho esfuerzos para preservar el violín de esta manera, y que han abandonado tales esfuerzos debido al daño al sonido, les digo: "Perseveren". Hagan tanto en tales esfuerzos como yo lo he hecho en diez años. Entonces, si aún no están satisfechos con los resultados, posiblemente mi método para preservar esas superficies internas de la desintegración les interese. Si me lo piden, ahora me considero ampliamente provisto de pruebas de que las superficies internas del violín pueden protegerse indefinidamente de la desintegración, y sin daño al sonido. Si mi interlocutor no tiene demasiados prejuicios, espero convencerlo de que el sonido de tales violines protegidos permanece igualmente dulce, mientras que las cualidades tonales de brillo e intensidad aumentan enormemente. Hubo una vez un hombre cuyo genio como imitador del violín Stradivarius, no solo llenó el mundo del violín con asombro, pero también, llenó los bolsillos de este hombre. Este hombre, JB Vuillaume, impostor-
Confesó que las superficies internas del violín no se pueden proteger sin dañar el sonido. Dado que este hombre era un impostor, cuando dijo "No", posiblemente quiso decir "Sí".
Resulta extraño cómo el "No" de un impostor sigue influyendo en el mundo del violín. En mi opinión, es una muestra de insensatez seguir el consejo de un estafador conocido. Creo que las declaraciones de quienes, incluso bajo juramento, no tienen otro objetivo en la vida que el dinero, deben tomarse con mucha cautela. De hecho, hay casos en los que vale más la pena. En este último grupo, me inclino por JB Vuillaume.
Como miembro del personal de un sanatorio de violines, tuve la oportunidad de presenciar el lado más lamentable de la desintegración de la superficie interior de un violín. A decir verdad, ver tal destrucción del valor del violín se convirtió en el agente estimulante para mi posterior esfuerzo por encontrar un medio seguro y un método para prevenirla. A menudo he presenciado la ruina total del valor tonal sin otra razón que dejar las superficies interiores del violín desprotegidas de la desintegración por la acción del calor y la humedad. La pregunta "¿Se puede prevenir tal desintegración sin dañar el sonido?" se presentó ante mis ojos en letras grandes hasta que actué en contra de una conclusión universal. Ahora, me sorprende que el daño al sonido, debido a la protección de las superficies interiores del violín, sea solo un temor imaginario. No afirmo que el sonido no pueda dañarse por la protección interior. Al contrario, sé muy bien que...
Que la protección de la superficie interior puede ser la causa de la pérdida de calidad tonal. Sé perfectamente que los métodos y medios rudimentarios para aplicarlos a las superficies internas del violín pueden dañar el sonido cuando se aplican a las superficies externas; pero no más ni menos. Sin embargo, al considerar el daño al sonido, no se debe pasar por alto lo siguiente: lo que daña el sonido cuando se aplica a una sola superficie de la caja de resonancia, multiplica dicho daño por el doble cuando se aplica a ambas superficies. En resumen, no se debe utilizar en el violín ningún medio ni método para proteger ninguna de las superficies, que resulte en un daño al sonido. Como músicos experimentados, saben bien que no es tanto lo que se toca, sino cómo se toca, lo que hace merecedor de un bis.
Si me hubiera rendido tras los primeros intentos de proteger las superficies internas de los violines, probablemente seguiría contemplando este problema desde la distancia. Si hubiera abandonado el esfuerzo, jamás habría recibido la aprobación de los 60 propietarios de violines así protegidos. Si el sonido de dichos violines se hubiera visto afectado por el tratamiento, no habría recibido la aprobación de los propietarios. En todos los casos, mis honorarios estaban sujetos a aprobación. Nunca los he perdido. Menciono estos últimos hechos únicamente como prueba de que la protección de las superficies internas del violín no perjudica su sonido.
[Quizás sea apropiado afirmar que trabajé en violines durante más de cuarenta años sin ninguna intención ni expectativa alguna de convertir ese trabajo en un negocio.]

Solo cobraba una tarifa cuando me llegaban violines en tal cantidad que requerían mucho de mi tiempo.
En mi opinión, aplicar protección tanto a las superficies exteriores como interiores de las tapas de un violín, sin dañar el sonido, exige experiencia y criterio que solo se adquieren con la práctica.
«Doctor, ¿cuánto tiempo durarán sus violines protegidos?» No puedo responder a esa pregunta con más detalle que decir que no veo ninguna razón por la que el material que utilizo para la protección de la superficie interior no dure igual que la protección de la superficie exterior. Para dicha protección de la superficie interior, utilizo goma copal como base, pero la dureza del copal se reduce con gomas mucho más suaves, resistentes y elásticas como el mástique, el elemí, el sandáraca, etc.
[Los detalles de este trabajo se darán a conocer más adelante.]
Nuevamente les llamo la atención sobre el hecho de que la madera de la caja de resonancia, que produce la mejor calidad de sonido, es la primera en arruinarse por la desintegración causada por el calor y la humedad. Sé cuánta variedad existe en cuanto al gusto por el violín. Recuerdo bien mis propios gustos diferentes en distintos períodos de mi experiencia. En los primeros períodos, no exigía nada más que una gran potencia en el sonido del violín. La calidad, ya sea en tonos simples o dobles, la sacrificaba voluntariamente a la mera potencia del sonido. Al mirar hacia atrás a mis primeros intentos de hacer temblar la tierra con mi gran sonido, me asombra la montaña de sufrimiento que vertí con entusiasmo en el oído de mis amigos cercanos. La palabra "cercanos" solo se refiere a la proximidad, y por
De ninguna manera se refiere a "amar" a los amigos. Hoy, cuando hablo de la caja de resonancia que produce el mejor tono, me refiero a la mayor belleza tonal, o al tono más agradable, o al tono más dulce; y me refiero especialmente a la gran belleza de los tonos de doble cuerda.
Como muestra del deterioro que el calor y la humedad provocan en la madera de mejor calidad de la caja de resonancia, presento este violín antiguo. Se desconoce el fabricante. Su aspecto desgastado hace innecesario cualquier certificado de antigüedad. ¡Los orificios de las clavijas están desgastados más allá del tamaño de cualquier clavija de violín! La empuñadura está profundamente desgastada por el roce del pulgar y el índice. En las superficies tocadas por la barbilla y el hombro, el barniz ha desaparecido. En los extremos superior e inferior de la caja de resonancia, el tejido conectivo celular blando, entre la parte más densa de la veta, se ha desgastado, dejando la parte más densa en forma de crestas. Debido a su gran influencia en los tonos de doble cuerda, hago especial hincapié en la blandura del tejido conectivo en esta muestra de madera de la caja de resonancia. Gracias a mi dilatada experiencia, puedo afirmar con seguridad que los tonos de doble cuerda de este violín poseían una belleza notable.
[Utilizo el tiempo pasado porque este violín no está en condiciones de ser tocado.]
La única afirmación en los libros de texto de filosofía en la que he encontrado algo valioso para la selección de la madera de la caja de resonancia del violín, está contenida en la siguiente frase: "La caja de resonancia puede compararse con un haz de cuerdas".
Considero que esta afirmación es Evidente- valiosa.

En los libros de texto, la palabra "cuerdas" se refiere a la parte más densa del grano. Para que sean útiles en la caja de resonancia, las "cuerdas" deben mantenerse unidas por un medio de conexión. Evidentemente, la naturaleza de dicho medio de conexión modifica la acción de esas cuerdas. Por lo tanto: la libertad de acción de esas "cuerdas" es como la rigidez del tejido conectivo. Como se demostró anteriormente con un ejemplo, la caja de resonancia rígida del violín no puede producir tonos dobles agradables, porque la rigidez del tejido conectivo interfiere con la acción independiente de las fibras contiguas ("cuerdas", de los libros de texto). En esta antigua caja de resonancia, el tejido conectivo es blando.
El dueño de este violín tan desgastado es el Sr. August Wolfe, Director Musical del Valparaiso College, Valparaiso, Indiana. El profesor Wolfe trajo este viejo violín de su casa en Austria. Cuando era niño, este viejo instrumento desgastado le llegó como regalo de su padre, acompañado del comentario falaz habitual: "Servirá para empezar". [Considerando el estado actual de este violín, creo que solo un niño de habla alemana podría haber sobrevivido para convertirse en un maestro.
violinista.]
En tamaño, este violín se conoce como i. Les pido que observen la belleza extremadamente delicada de este antiguo mástil y cabeza. Ni siquiera el de Hebe es más hermoso. Esta mano esbelta, estas paredes delgadas del clavijero, esas exquisitas líneas de estrías y voluta, apelan poderosamente a nuestro sentido de la belleza. Debido a que los agujeros de las clavijas están desgastados, y una fractura se extiende a través del agujero de la clavija A, y una pieza
Se ha perdido el borde del acanalado, el profesor sugiere reemplazar este cuello viejo y desgastado por uno nuevo. Hasta ahora he realizado trabajos quirúrgicos que requieren nervios, pero, al separar esta hermosa cabeza de su cuerpo, me detengo. "Es un caso de insuficiencia cardíaca". Podría dedicar fácilmente días al trabajo de restaurar las líneas rotas en este cuello y cabeza de aspecto hebeano para que esas líneas de belleza permanezcan para deleitar la vista del conocedor.
Ahora voy a presentar algo inusual. Quiero llamar la atención sobre el barniz de este viejo violín. Este barniz es blando; demasiado blando para los estándares modernos. Además, a juzgar por la cantidad que se encuentra ahora en las cavidades, la cantidad de barniz aplicada originalmente era mucho mayor que la que se usa hoy en día. Es tan blando que, con una presión moderada, mi dedo deja una marca perceptible. Dado que la fricción produce el olor característico del mástique, supongo que esta goma blanda está en exceso. En este barniz, me interesan menos las gomas que el colorante empleado. Como pueden observar, los colores predominantes son el rojo y el negro, mezclados en cantidades que dan como resultado un tono marrón rojizo intenso.
¡Pero qué color marrón rojizo tan sucio! Ningún luthier moderno que se precie permitiría que un trabajo de color tan sucio y de aspecto fangoso surgiera de su montón de madera de desecho.
"Doctor, ¿no es este un violín antiguo muy raro?" ¡Claro!
[Avisé que iba a presentar algo inusual.]
Ahora voy a presentar algo aún más sorprendente: 
más impactante y más patético. Voy a presentar un ejemplo de desintegración de la madera de la caja de resonancia que quedará grabado en su memoria; un ejemplo de combustión lenta de la madera en presencia de calor y humedad; un ejemplo que muestra el camino cuesta abajo recorrido por esas invaluables joyas antiguas del arte del sonido, ahora varadas en el último siglo de la historia del violín. Esta pérdida irreparable, y esta escena patética, podrían haberse evitado fácilmente y con certeza, estoy completamente convencido. Con una espátula delgada, retiro fácilmente esta vieja caja de resonancia.
¡A la sombra de Dante!
Como aplicación en el Hades, bien podríamos temer que la combustión lenta encabece la lista. La veracidad de aquel viejo dicho «El diablo está en el violín» ahora se confirma. Pero, extrañamente e inesperadamente, esta evidencia confirmatoria apunta al hecho de que «está en el violín» con el propósito de destruir el violín en lugar de destruir a los seres humanos. Se afirma que recientes descubrimientos arqueológicos sacan a la luz pruebas de algún error en las lecturas bíblicas. La evidencia proporcionada por este ejemplo de combustión lenta, complementada con otros ejemplos de naturaleza similar, podría ser suficiente para cambiar las lecturas ortodoxas de la siguiente manera: «El diablo es enemigo del violín».
El "punk" puede describirse como un producto de la madera por combustión lenta en presencia de calor y humedad. En la superficie interior de la caja de resonancia tenemos un ejemplo de "punk". Cae como polvo mezclado con finos hilos de fibra de madera cuando lo toco con la espátula. Esta condición de punk se extiende
Cubre toda la superficie interior de la caja de resonancia, y su profundidad es de 1,32 pulgadas. En algunos puntos, los revestimientos del borde superior de las costillas cuelgan en jirones parcialmente descompuestos. Se observa que tanto los revestimientos como los bloques de las esquinas son mucho más ligeros que los actuales.
Desconozco por qué no todas las muestras de madera para cajas de resonancia de la misma edad presentan condiciones de combustión similares. Solo puedo afirmar que las distintas muestras de madera poseen diferentes grados de resistencia a la combustión lenta.
[En mi juventud, la pólvora, el pedernal y el acero eran nuestros únicos fósforos. Tomábamos las mismas precauciones para mantener secos tanto la pólvora como el peltre. Así conservado, el peltre duraba mucho tiempo; pero, expuesto al calor y la humedad, como cuando se encontraba en el bosque, pronto se convertía en polvo. No todos los troncos en descomposición producen peltre de buena calidad. El buscador de peltre a menudo debe realizar una búsqueda prolongada antes de encontrar el mejor. En aquellos tiempos lejanos, según recuerdo, el peltre de mejor calidad se encontraba en ramas grandes y muertas de árboles vivos; o dentro del tronco de árboles que morían en pie. En la carpintería es bien sabido que la mayor profundidad de los cambios de color se encuentra en los árboles que mueren en pie; también, que la madera de los árboles que mueren así se descompone más rápidamente. Al buscar razones que expliquen la rapidez de desintegración observable en la mejor calidad de los violines antiguos, algunos autores sugieren que sus tapas armónicas se obtenían de árboles que habían muerto en pie. Sobre este punto no parece haber evidencia positiva. El hecho de que algunos
El hecho de que algunas muestras de cajas de resonancia se hayan deteriorado por la acción del calor y la humedad, mientras que otras permanecen intactas, sigue sin explicación.
Ahora sostengo este violín de tal manera que puedan mirar hacia la superficie interior de la tapa trasera. Hago hincapié en "mirar hacia" porque no pueden "mirar" esta tapa. Su mirada no puede penetrar este depósito de materia terrosa. El color de este depósito es negro, lo que sugiere aluvión. Pero, dado que el aluvión es un depósito de agua, debemos buscar en otra parte el origen de este depósito de polvo aéreo del "período Cremona".
[Algunos de ustedes pueden tener una opinión diferente. Admito que su opinión es tan buena como la mía, posiblemente mejor. Solo afirmo que soy incapaz de estimar el "polvo de Cremona" más allá de su valor nominal. Sé que algunos expertos en violines antiguos, "solo por su salud", (!), afirman que ningún otro polvo tiene
valor. No estoy "buscando la salud". Admito que el polvo de Cremona tiene tendencia a depositarse dentro del violín. De hecho, facilitar que dicho polvo se asiente ha mejorado la salud, (!), de muchos expertos. La buena salud es un bien preciado. Solo por esta razón, el "experto" en violines antiguos sin duda seguirá estando "fuera de juego".
Si intentara desempeñar el papel de "experto" en violines antiguos, podría continuar así: "Si se vieran algunas rocas dispersas en este depósito, indicarían que la edad de este viejo violín es igual a la época glacial. Pero, como no hay rocas, podemos continuar nuestra marcha con seguridad, en un
al revés. Sin duda, este depósito alguna vez fueron partículas voladoras de materia terrestre, polvo y nebulosidad. Ahora lo tenemos. El período nebuloso es el
límite.
Caballeros: ¡He aquí el primer violín!
Tras haber retirado suficiente punk de la caja de resonancia y suficiente suciedad de la placa posterior, podemos observar que ambas placas se cortan de la misma pieza. Este método de fabricación de las placas de violín parece haber prevalecido más en los inicios de la historia del violín que en la actualidad. En mi experiencia, la placa entera produce un sonido igual de bueno que la placa dividida, siempre que la veta de la caja de resonancia forme un ángulo recto, o casi recto, con la curvatura en la zona de producción del sonido. Sin embargo, cuando dicho ángulo de la veta es oblicuo a la curvatura, siempre he notado una pérdida de potencia y brillo en el sonido. Esta misma pérdida se produce también en la caja de resonancia dividida cuando el ángulo de la veta es oblicuo a la curvatura. Atribuyo esta pérdida a una menor acción de las fibras.
[Honeyman afirma con positividad que siempre que se encuentra una caja de resonancia de Nicolas Amatus con espesores menores a 1 en el centro y 1 en los bordes, dicha caja de resonancia ha sido sometida a una regradación por algún artesano moderno. Por el bien de la música y de la humanidad, considero que dicha regradación es afortunada. Mi experiencia con cajas de resonancia tan gruesas y de densidad media es que incluso un juego de cuerdas de calibre 1 no puede superar la rigidez lo suficiente como para producir una octava de
[Buen tono en cada cuerda.] Como evidencia que respalda la confusión encontrada en la historia del violín, cito aquí al francés N.E. Simoutre, libro y gráficos, París, 1885. Simoutre presenta gráficos de dos violines Nicolas Amati sin fecha: Grosor, Lámina 1, centro de la caja de resonancia, 360 mm (es decir, tres y cinco décimas de milímetro) hasta 250 mm, y en los bordes, 400 mm (es decir, 4 milímetros). Grosor, Lámina 2, en el centro 450 mm, hasta 200 mm, y en los bordes, 300 mm. (Simoutre explica sus cifras por centième des millimeters, que significa tantas centésimas de milímetro).
milímetro. )
Al comparar el grosor en el centro de la caja de resonancia de Nicolas Amatus, según lo indicado por Honeyman, con el grosor máximo proporcionado por Simoutre, encontramos una diferencia de siete centésimas de pulgada.
Considerando la profunda devoción humana por el violín, me asombra el maltrato que ha sufrido. Si este viejo violín fue en su día una obra maestra del arte sonoro, el abuso al que ha sido sometido es suficiente para arruinarlo. Reparar este "viejo instrumento desgastado" para prolongar su vida útil supone una dura prueba para la paciencia. Frágil como el cristal, requiere un trato cuidadoso. En este trabajo, el dinero importa poco. La recompensa reside en escuchar unos tonos tan dulces que el dinero no puede reproducir.
¡Dulce violín antiguo! Desgastado,
Desgarrado, raspado desde el amanecer hasta el amanecer. El trabajo para ti siempre ha sido repartir alegría desde que naciste. La senilidad, ahora dentro de tu forma,
Te considera una "cosa vieja y desgastada"; sin embargo, tu tono hacia nosotros permanece
Una joya inalcanzable para cualquier rey.

